%% ma: L10-B14-0032
%% lop: 10
%% chuong: 5
%% bai: 14
%% dang: 10.14.5
%% loai: DS
%% muc: [NB, TH, TH, VD]
%% dap_an: "ĐSSS"
%% nguon: "chuyen-de-cac-so-dac-trung-cua-mau-so-lieu-khong-ghep-nhom-toan-10-chuong-trinh-moi.pdf (GV. Phan Nhật Linh) (P04), Bài 03 Dạng 1 Phần II Câu 9"
%% kien_thuc: "Lớp 10 Bài 13 (số trung bình, tứ phân vị), Bài 14 (khoảng biến thiên, khoảng tứ phân vị, giá trị bất thường)"
%% trang_thai: chinh-thuc
%% ghi_chu: "Đáp án gốc ý a) Sai là sai: lời giải gốc lấy giá trị cuối dãy chưa sắp xếp (24 - 21 = 3); đúng R = 30 - 21 = 9 nên a) Đúng. Tự duyệt theo yêu cầu giáo viên 22/10/2026."
\begin{ex}
Nhiệt độ trung bình (đơn vị: $^\circ$C) theo tháng tại một thành phố trong năm 2023 được cho như sau:
\begin{center}
\begin{tabular}{|c|c|c|c|c|c|c|c|c|c|c|c|}\hline
$21$ & $23$ & $25$ & $27$ & $28$ & $30$ & $26$ & $28$ & $25$ & $23$ & $23$ & $24$\\\hline
\end{tabular}
\end{center}
Xét tính đúng sai của các khẳng định sau:
\choiceTF{\True Khoảng biến thiên của mẫu số liệu là $9$}{Nhiệt độ trung bình năm 2023 của thành phố là $25\,^\circ$C}{Tứ phân vị thứ ba của mẫu số liệu là $24$}{Giá trị bất thường của mẫu số liệu là $30$}
\loigiai{Sắp xếp: $21$; $23$; $23$; $23$; $24$; $25$; $25$; $26$; $27$; $28$; $28$; $30$.

a) Đúng: $R=30-21=9$.

b) Sai: $\overline x=\dfrac{303}{12}=25{,}25$.

c) Sai: $Q_3=\dfrac{27+28}2=27{,}5$.

d) Sai: $Q_1=23$, $\Delta_Q=4{,}5$, $Q_3+1{,}5\Delta_Q=34{,}25>30$ nên $30$ không phải giá trị bất thường.}
\end{ex}
